\scanpage{287}
% archive sheet number: 288
(1) $S$ anneau \inserted{$\mathbb Z$-}gradué \inserted{\unsure{comm|comm.}}, \uline{à degrés $\geq0$}, $\displaystyle S=\sum_{n\geq0} S_{\inserted{n}}^{\deleted{\unsure{n}}}$

\deleted{$\unsure{S_0=k}$}

(2) $k=S_{\inserted{0}}^{\deleted{\unsure{0}}}$

(3) $\operatorname{Modgrad}(S)=\mathcal M$

(4) $\underline t$ foncteur translation
\[
\left\{\begin{aligned}
&\underline t:\mathcal M\mathrel{\unsure{-|\to}}\mathcal M\\
(4.1)\quad &\underline t(M)_i=M_{i+1}
\end{aligned}\right.
\qquad\text{c'est un \uline{automorphisme} de $\mathcal M$}
\]

(4.2) \deleted{\unsure{M|$\mathcal M$}} \unsure{on pose} \inserted{$M[i]=\underline t^{\,i}(M)$.}

Pour tout entier $N$, on considère les deux sous-catégories \inserted{\unsure{str.|strict.}} pleines \inserted{$\mathcal M_N,\mathcal M^{N}$} de $\mathcal M$ \unsure{définies} par
\[
\left\{\begin{aligned}
\operatorname{Ob}\mathcal M_N&=\{P\in\operatorname{Ob}\mathcal M\mid P_i=0\quad\forall i<N\}\\
\operatorname{Ob}\mathcal M^{N}&=\{P\in\operatorname{ob}\mathcal M\mid P_i=0\quad\forall i\unsure{>|{\geq}}N\}
\end{aligned}\right.\tag*{(5)}
\]

\uline{NB} Les $\mathcal M_N$ ($\mathcal M^{N}$) décroissent (croissent) pour $N$ croissant. On a
\[
\left\{\begin{aligned}
t^{i}(\mathcal M_N)&=\mathcal M_{N-i}\\
t^{i}(\mathcal M^{N})&=\mathcal M_{\deleted{\unsure{N-i}}}^{N-i}
\end{aligned}\right.\tag*{(6)}
\]
\deleted{\unsure{donc}}

En particulier \deleted{\unsure{(6)|est}}
\[
\left\{\begin{aligned}
\deleted{\unsure{t}}\,\mathcal M_N&=t^{-N}\mathcal M_0\quad,\\
\mathcal M^{N}&=t^{-N}\mathcal M^{0}
\end{aligned}\right.\tag*{(6.1)}
\]

\inserted{Dans} la catégorie $\mathcal M$, munie déjà de l'\unsure{autom.} $t$, la donnée des familles $(\mathcal M_N)$ et $(\mathcal M^N)$ avec leurs inclusions définitoires
\[
\mathcal M_N\supset\mathcal M_{N'}\quad\text{et}\quad \mathcal M^N\subset\mathcal M^{N'}\quad\text{si}\quad N\leq N'\tag*{(7)}
\]
équivaut à celle \deleted{\unsure{des}} du couple
\[
\mathcal M_0,\mathcal M^0\quad\text{avec}\quad t\mathcal M_0\unsure{\subset|\supset}\mathcal M_0,\quad t\mathcal M_{\deleted{\unsure{0}}}^{0}\supset\mathcal M_{\deleted{\unsure{0}}}^{0}\tag*{(8.)}
\]
On va voir d'ailleurs comment $\mathcal M_{\unsure{N}},\mathcal M^{N-1}$ \deleted{(\unsure{ou encore} $\unsure{\mathcal M_N,\mathcal M^N}$)} se déterminent mutuellement -- donc $\mathcal M_0$ \uline{ou} $\mathcal M^0$ permet de \unsure{retrouver|retirer} tous \deleted{\unsure{les}} $\mathcal M_N,\mathcal M^N$.
